\section{Introduction}
\label{sec:introduction}

In the number 98, every French child learns to say ``quatre-vingt-dix-huit''---literally, ``four-twenty-ten-eight,'' encoding the arithmetic expression $4\times20+10+8$. This vigesimal (base-20) counting system, inherited from Gaulish and Celtic linguistic influence, is one of the most notoriously counterintuitive features of French. Meanwhile, Belgian and Swiss French speakers say simply ``nonante-huit'' (ninety-eight), using a transparent decimal structure. This linguistic divergence creates a natural experiment: when a language model processes these two surface forms, does it arrive at the same internal numeric representation?

Recent work has revealed that large language models encode surprisingly precise base-10 digit representations in their hidden states \citep{levy2024language, kadlcik2025pretrained}. \citet{levy2024language} showed that circular probes can recover individual digits from number token representations with over 90\% accuracy, and that these representations causally influence model outputs. However, this line of work has focused almost exclusively on digit strings and simple English word forms. Whether models can decompose the complex arithmetic structure of vigesimal French number words into clean base-10 representations remains unknown.

{\bf Why does this matter?} Understanding how models process numbers embedded in complex linguistic structures reveals fundamental aspects of the interface between language and mathematics in transformers. If a model can internally ``translate'' ``quatre-vingt-dix-huit'' ($4\times20+10+8$) into the digit representation $[9, 8]$, it demonstrates that distributional learning has captured not just word co-occurrences but implicit arithmetic structure. Conversely, if vigesimal forms produce degraded representations, it points to fundamental limits of how linguistic structure constrains numerical reasoning.

We apply circular digit probes \citep{levy2024language} to \mistral across four input formats: digit strings, English number words, France French number words (vigesimal), and Belgian French number words (decimal). Our probes operate on hidden representations extracted at the last token position and predict each base-10 digit via circular (cosine/sine) encodings. We evaluate on 1,000 numbers (0--999) with a stratified train/test split.

Our results reveal three key findings. First, French number word representations---including vigesimal forms---encode base-10 digits at 98.0\% accuracy, comparable to English words (98.5\%) and significantly outperforming digit strings (91.0\%, $p=0.002$ by McNemar's test). Second, the vigesimal structure causes only localized difficulty in the \emph{soixante-dix} (70s) range (88.2\%), with errors revealing partial encoding of the ``soixante'' (60) component. Third, number word representations peak at early layers (5--6), while digit string representations peak at the final layer (31), suggesting fundamentally different processing pathways.

We make the following contributions:
\begin{itemize}[leftmargin=*,itemsep=0pt,topsep=0pt]
    \item We conduct the first systematic study of how LLMs represent French vigesimal numbers internally, showing that base-10 digit information is recoverable at 98.0\% accuracy despite the complex arithmetic surface form.
    \item We exploit the natural experiment between France French (vigesimal) and Belgian French (decimal) to isolate the effect of counting system structure on internal representations, finding that Belgian French achieves perfect accuracy on vigesimal-range numbers where France French drops to 88.2\%.
    \item We discover that number word representations emerge early in the network (layers 5--6) while digit string representations peak late (layer 31), revealing distinct processing pathways for semantic versus tokenized numerical inputs.
    \item We provide a detailed error analysis showing that vigesimal errors reflect partial arithmetic decomposition---the model sometimes encodes ``soixante'' as 60 rather than the composed value 70+---offering a window into the model's internal processing of compositional number words.
\end{itemize}

The remainder of this paper is organized as follows. \Secref{sec:related_work} reviews related work on number representations in LLMs. \Secref{sec:methodology} describes our probing methodology and experimental setup. \Secref{sec:results} presents our main results, layer-wise analysis, and error analysis. \Secref{sec:discussion} discusses implications and limitations. \Secref{sec:conclusion} concludes.
